% !TeX program = xelatex
\documentclass[UTF8,a4paper,12pt]{ctexart}
\usepackage[a4paper,margin=2.5cm]{geometry}
\usepackage{booktabs}
\usepackage[colorlinks=true,citecolor=blue]{hyperref}

\title{面向复杂正常背景的开放集焊缝缺陷实例定位}
\author{作者：王俊杰\\单位：北京理工大学}
\date{\today}

\begin{document}
\maketitle

\begin{abstract}
\noindent\textbf{问题：}完整焊缝射线图中的复杂正常结构可能造成误报，闭集模型难以处理训练时未见的缺陷类别。

\noindent\textbf{数据：}拟使用可追溯的焊缝射线图像；具体范围和标注情况待审计。

\noindent\textbf{方法：}拟研究类别无关候选生成、背景条件化正常建模和已知/未知拒识。

\noindent\textbf{预期指标：}未知缺陷实例召回、误报/原片、已知类检测性能和候选覆盖率。

\noindent\textbf{当前状态：}研究计划骨架，尚无真实实验结果。
\end{abstract}

\noindent\textbf{关键词：}焊缝检测；开放集识别；异常定位；复杂背景

\section{引言}
介绍研究背景、问题定义、研究目标和可检验假设。

\section{相关工作}
综述焊缝射线检测、工业异常定位和开放集检测。
PatchCore 使用正常图像的局部特征记忆库进行异常检测与定位
\cite{roth2022patchcore}。

\section{数据与实验协议}
说明数据来源、标签映射、按焊缝分组的训练/验证/测试划分和防泄漏措施。

\section{方法}
描述类别无关候选、区域特征、背景条件建模和未知拒识。

\begin{figure}[htbp]
\centering
\fbox{\parbox[c][3.5cm][c]{0.85\linewidth}{
\centering
完整射线原片 $\rightarrow$ 类别无关候选 $\rightarrow$ 区域特征\\[0.6em]
背景条件 $\rightarrow$ 正常/异常评分 $\rightarrow$ 已知/未知判断\\[0.6em]
原片坐标还原与实例评估
}}
\caption{拟议方法流程占位图，后续替换为正式图。}
\label{fig:framework}
\end{figure}

\section{实验设计}
说明对比方法、候选预算、阈值校准和统计分析。

\section{结果}
真实实验完成后填写，不要把模拟结果写成实测结果。

\begin{table}[htbp]
\centering
\caption{主要结果占位}
\begin{tabular}{lccc}
\toprule
方法 & 未知实例召回 & 误报/原片 & 已知类 mAP \\
\midrule
全局正常基线 & 待实验 & 待实验 & 待实验 \\
背景条件方法 & 待实验 & 待实验 & 待实验 \\
\bottomrule
\end{tabular}
\end{table}

\section{讨论}
分析结果、失败案例、数据限制和适用范围。

\section{结论}
根据真实实验结果总结工作和结论。

\bibliographystyle{plain}
\bibliography{references}

\end{document}
